\chapter{Shell Commands and Pipes}
\label{ch:shell}

\index{shell commands}
\index{pipe commands}

One of acme's most powerful features is its seamless integration with
external commands.  Any text can be executed as a shell command.  The
Unix toolkit becomes your extension mechanism.

\section{Running External Commands}
\label{sec:external}

\index{executing commands!external}

When you middle-click (\button{2}) on text that is not a built-in
command, acme runs it as a shell command.  The shell used is determined
by the \env{acmeshell} environment variable, defaulting to \cmd{sh}.

The command runs in the directory associated with the current window.
For a file \file{/home/user/project/main.c}, commands execute in
\file{/home/user/project/}.

\subsection{Environment}

Acme sets the following environment variables for child processes:

\begin{itemize}
\item \env{winid} --- the numeric ID of the window.
\item \env{\%} --- the filename of the window.
\item \env{samfile} --- same as \env{\%}.
\end{itemize}

\subsection{Output}

Standard output and standard error from the command appear in the
\file{+Errors} window for that directory.  If no \file{+Errors} window
exists, one is created automatically.

\subsection{Examples}

Type any of these in a tag bar and middle-click to execute:

\begin{lstlisting}[language=bash]
date                     # show the date
ls -la                   # list files
make                     # build the project
grep -rn TODO *.c        # search for TODOs
git status               # check version control
\end{lstlisting}

\section{Pipe Prefixes: <, >, |}
\label{sec:pipes}

\index{pipe!< prefix}
\index{pipe!> prefix}
\index{pipe!| prefix}

Commands prefixed with \texttt{<}, \texttt{>}, or \texttt{|} interact
with the current selection:

\begin{longtable}{@{}cp{3.5in}@{}}
\toprule
Prefix & Behavior \\
\midrule
\endhead
\texttt{<cmd} & Run \texttt{cmd}, \textbf{replace the selection}
  with its standard output. \\
\texttt{>cmd} & Run \texttt{cmd}, sending the \textbf{selection as
  standard input}.  The selection is unchanged. \\
\texttt{|cmd} & Run \texttt{cmd}, sending the \textbf{selection as
  standard input} and \textbf{replacing it} with standard output. \\
\bottomrule
\end{longtable}

\subsection{Examples}

\paragraph{Replace selection with command output (\texttt{<}):}

\begin{lstlisting}[language=bash]
<date              # insert the current date
<seq 1 10          # insert numbers 1 through 10
<cat /etc/hostname # insert the hostname
\end{lstlisting}

\paragraph{Send selection to command (\texttt{>}):}

\begin{lstlisting}[language=bash]
>wc            # count words/lines in selection
>lp            # print selection
>xclip         # copy to X clipboard
>mail user@x   # email the selection
\end{lstlisting}

\paragraph{Filter selection through command (\texttt{\textbar}):}

\begin{lstlisting}[language=bash]
|sort             # sort selected lines
|sort -u          # sort and remove duplicates
|fmt -w 72        # reflow text to 72 columns
|tr a-z A-Z       # uppercase the selection
|awk '{print NR, $0}'  # add line numbers
|sed 's/old/new/g'     # search and replace
|column -t        # align columns
|indent           # reformat C code
|gofmt            # reformat Go code
\end{lstlisting}

These pipe commands are the primary way to extend acme.  Instead of
plugins, you use the standard Unix toolkit.

\section{The +Errors Window}
\label{sec:errors}

\index{+Errors window}

Command output goes to a window named with the directory path plus
\file{+Errors}.  For example, commands run from
\file{/home/user/project/main.c} produce output in
\file{/home/user/project/+Errors}.

The \file{+Errors} window is an ordinary acme window.  You can:

\begin{itemize}
\item Search it with \button{3}.
\item Right-click on error messages like \file{main.c:42: error} to
      jump directly to the source location.
\item Close it with \cmd{Del}.
\item Clear it by selecting all and deleting.
\end{itemize}

\section{The Win Command}
\label{sec:win}

\index{Win}
\index{PTY}
\index{interactive shell}

The \cmd{Win} command opens an interactive shell session inside an acme
window.  It creates a pseudo-terminal (PTY) connected to your shell
(typically bash or whatever \env{SHELL} is set to).

\subsection{How It Works}

\begin{enumerate}
\item Execute \cmd{Win} (middle-click it in any tag bar).
\item A new window opens, named after the current directory with a
      \file{-win} suffix.
\item Type commands and press \key{Enter} to send them to the shell.
\item Output from the shell appears in the window body.
\end{enumerate}

\subsection{Behavior}

The Win window is not a terminal emulator---it is an acme text buffer
connected to a shell via PTY.  This means:

\begin{itemize}
\item You can edit text before pressing \key{Enter}.  The shell does not
      see your keystrokes until you press \key{Enter}.
\item You can select output text and use it like any other acme text
      (right-click filenames, middle-click commands, etc.).
\item ANSI escape sequences are stripped.  Colors and cursor positioning
      do not work.
\item The PTY is set to single-column width, so \cmd{ls} and similar
      commands produce one-entry-per-line output.
\item Echo is disabled---acme shows what you type, not the terminal.
\end{itemize}

\subsection{Shell Configuration}

The Win command starts the shell with a minimal environment:

\begin{itemize}
\item \texttt{TERM=dumb} --- suppresses escape sequences.
\item \texttt{PS1="\$ "} --- a simple prompt.
\item \texttt{PROMPT\_COMMAND}, \texttt{COLORTERM}, and
      \texttt{LS\_COLORS} are unset.
\end{itemize}

The shell is started with \cmd{--noediting --norc --noprofile -i}
flags (for bash) to minimize interference.

\subsection{Closing}

Close a Win window with \cmd{Del} like any other window.  The shell
process receives SIGHUP and terminates.

\section{When to Use What}

\begin{longtable}{@{}lp{3.5in}@{}}
\toprule
Need & Approach \\
\midrule
\endhead
Run a one-off command & Type it in the tag, \button{2} click \\
Transform selected text & Use \texttt{|cmd} \\
Insert command output & Use \texttt{<cmd} \\
Send text to a command & Use \texttt{>cmd} \\
Interactive session & Use \cmd{Win} \\
Navigate compiler errors & Right-click error lines in \file{+Errors} \\
\bottomrule
\end{longtable}
